\section{n3. An intrinsic incidence criterion and its precise remaining case}
\textit{Supports: ENTIRE qualified G331; own n1; own n2.}

Use n2's general excellent reduced setting. Define the finite graph on the
actual normal LOCAL CENTRES by an edge $i\ne j$ exactly when
$\rho(Y_i)\cap\rho(Y_j)\ne\varnothing$. Equivalently, by G331's exact
infinite-conductor rule, there are centred normalized semivaluations on
the two local conductor quotients restricting to the same semivaluation
on $R/\mathfrak c$. This uses EVERY prime support, including the infinite
stratum, rather than only generic valuations or irreducible components.
It is an intrinsic boundary criterion; no finite algorithm deciding all
of these semivaluation tests is claimed.

Write $\operatorname{CE}(R)$ for: every finite continuous signed obstacle
on its full primary section has a greatest geometric $C_R$ minorant.
By n1 the same predicate results from greatest $D_R$, raw continuity or
uniform every-degree canonical convergence. N2 proves the new necessary
condition
\[
                 \operatorname{CE}(R)\ \Longrightarrow\
                 \text{the incidence graph has no edges}.       \tag{11}
\]
There is no smoothness, characteristic or source-CE condition in (11).

There is also a complete scoped classification. Assume that restriction
$\rho|_{Y_i}$ is INJECTIVE for each single centre. This hypothesis is
equivalent to the local conductor quotient restriction being injective
on its entire normalized boundary; finite-conductor fibres are already
singletons. Then
\[
 \operatorname{CE}(R)\quad\Longleftrightarrow\quad
 \left[\text{no graph edges AND }\operatorname{CE}(S_i)
                  \text{ for every centre }i\right].             \tag{12}
\]
Proof: no edges plus the assumed internal injectivity makes $\rho$ a
continuous bijection between compact Hausdorff boundaries, hence a
homeomorphism. G331's full-fibre POSITIVE descent makes its pullback an
order isomorphism $C_R\simeq\prod_i C_{S_i}$, and likewise for $D$.
Thus an arbitrary continuous obstacle corresponds to an arbitrary tuple
of continuous obstacles, and both raw and attained envelopes correspond
componentwise. This proves both directions when there are no edges.
If there is an edge, n2 excludes $\operatorname{CE}(R)$, proving (12).
In particular, whenever these normal centres have CE, internal
injectivity leaves precisely the intrinsic graph test for target CE.
This theorem applies even when no finite regular cover of any centre is
supplied. It does not prove CE of the normal centres themselves.

The complement of this scope is precise. In general (11) still holds,
but no edges does not decide identifications WITHIN a single centre.
For those internal conductor quotient fibres the present argument does
not furnish the zero-versus-negative centre separation used in n2.
General normal higher-dimensional nonisolated CE, higher-dimensional
positive/mixed-characteristic regular-source CE, and arbitrary
same-centre conductor classification remain unresolved. Consequently
this is a new complete scoped criterion and an original-residual bridge,
not a complete solution of the unchanged P12.3.

For another primary ideal $\mathfrak b$, put $d(v)=v(\mathfrak b)>0$.
The granted homogeneous transport is
$f_{\mathfrak b}(v/d(v))=f_{\mathfrak a}(v)/d(v)$, with the same
denominator on every full normalization fibre. The order bijection of
competitors gives the same formula for raw or attained envelopes.
Since $d,d^{-1}$ are bounded continuous, CE, graph edges and internal
injectivity are gauge independent. N2's jump becomes
$\kappa/d(v_*)$, because $d(v_\lambda)\to d(v_*)$.
The granted exact completion homeomorphism and primary-ideal contraction
transport the obstruction and CE predicate to completion. This statement
does not assert equivalence of model categories or equality of generic
degrees after base change. Nilpotents first reduce away; dimension-zero
local reduction has the unique empty boundary and empty-envelope rule.

For completeness the curve comparison is reconstructed, not claimed new.
An excellent reduced curve has finitely many normal local DVR centres.
Their normalized sections are singletons, their conductor values finite,
and the full normalization restriction is a bijection. Primary powers,
unit profiles and positive real scaling yield all finite nonpositive
values at a DVR; actual decreasing constants yield all extended values.
Semilocal gluing and granted geometric positive descent therefore give
\[
                           P_R(f)=\min(f,0)
\]
on the finite branch link. This is a consequence of actual Cartier
comparison and descent. The graph has no edges here, consistent with
(11). Literal classical/CLD logarithmic radial extensions keep their
multibranch mixed-zero obstruction and are not identified with this
valuative boundary class. The qualified published smooth specialization
of the class remains exactly G331's; no source-CE enlargement is inferred.

Artifacts: n3.tex and P12.3-conductor-centre-obstruction.tex. Rerun:
python delivery.py. Equation (12) is a mathematical conditional
classification; matching artifact hashes and finite checks do not
verify its premises or establish independent mathematical acceptance.
